\Chapter{3}

\Verse{1}
{
    \zA{却说雨村忙回头看时，不是别人，乃是当日同僚一案参革的张如圭。}
}
{
    \eA{But to proceed with our narrative. Yü-ts'un, on speedily turning round,
        perceived that the speaker was no other than a certain Chang Ju-kuei, an old
        colleague of his, who had been denounced and deprived of office,}
}
{
}

\Verse{2}
{
    \zA{他系此地 人，革后家居，今打听得都中奏准起复旧员之信，他便四下里寻情找门路，忽遇见 雨村，故忙道喜。}
}
{
    \eA{on account of some case or other; a native of that district, who had, since
        his degradation, resided in his family home. Having lately come to hear the
        news that a memorial, presented in the capital, that the former officers
        (who had been cashiered) should be reinstated, had received the imperial
        consent,}
}
{
}

\Verse{3}
{
    \zA{二人见了礼，张如圭便将此信告知雨村，雨村欢喜，忙忙叙了两 句，各自别去回家。}
}
{
    \eA{he had promptly done all he could, in every nook and corner, to obtain
        influence, and to find the means (of righting his position,) when he,
        unexpectedly, came across Yü-ts'un, to whom he therefore lost no time in
        offering his congratulations. The two friends exchanged the conventional
        salutations,}
}
{
}

\Verse{4}
{
    \zA{冷子兴听得此言，便忙献计，令雨村央求林如海，转向都中去 央烦贾政。}
}
{
    \eA{and Chang Ju-kuei forthwith communicated the tidings to Yü-ts'un. Yü-ts'un
        was delighted, but after he had made a few remarks, in a great hurry, each
        took his leave and sped on his own way homewards. Leng Tzu-hsing, upon
        hearing this conversation, hastened at once to propose a plan,}
}
{
}

\Verse{5}
{
    \zA{雨村领其意而别，回至馆中，忙寻邸报看真确了，次日面谋之如海。}
}
{
    \eA{advising Yü-ts'un to request Lin Ju-hai, in his turn, to appeal in the
        capital to Mr. Chia Cheng for support. Yü-ts'un accepted the suggestion, and
        parted from his companion. On his return to his quarters,}
}
{
}

\Verse{6}
{
    \zA{如 海道：“天缘凑巧，因贱荆去世，都中家岳母念及小女无人依傍，前已遣了男女船 只来接，因小女未曾大痊，故尚未行，此刻正思送女进京。}
}
{
    \eA{he made all haste to lay his hand on the Metropolitan Gazette, and having
        ascertained that the news was authentic, he had on the next day a personal
        consultation with Ju-hai. "Providence and good fortune are both alike
        propitious!" exclaimed Ju-hai. "After the death of my wife, my
        mother-in-law, whose residence is in the capital, was so very solicitous on
        my daughter's account,}
}
{
}

\Verse{7}
{
    \zA{因向蒙教训之恩，未经 酬报，遇此机会岂有不尽心图报之理。}
}
{
    \eA{for having no one to depend upon, that she despatched, at an early period,
        boats with men and women servants to come and fetch her. But my child was at
        the time not quite over her illness,}
}
{
}

\Verse{8}
{
    \zA{弟已预筹之，修下荐书一封，托内兄务为周 全，方可稍尽弟之鄙诚；即有所费，弟于内家信中写明，不劳吾兄多虑。”}
}
{
    \eA{and that is why she has not yet started. I was, this very moment, cogitating
        to send my daughter to the capital. And in view of the obligation, under
        which I am to you for the instruction you have heretofore conferred upon
        her, remaining as yet unrequited, there is no reason why, when such an
        opportunity as this presents itself,}
}
{
}

\Verse{9}
{
    \zA{雨村一 面打恭，谢不释口，一面又问：“不知令亲大人现居何职？}
}
{
    \eA{I should not do my utmost to find means to make proper acknowledgment. I
        have already, in anticipation, given the matter my attention, and written a
        letter of recommendation to my brother-in-law,}
}
{
}

\Verse{10}
{
    \zA{只怕晚生草率，不敢进 谒。”}
}
{
    \eA{urging him to put everything right for you, in order that I may, to a
        certain extent,}
}
{
}

\Verse{11}
{
    \zA{如海笑道：“若论舍亲，与尊兄犹系一家，乃荣公之孙：大内兄现袭一等将 军之职，名赦，字恩侯；二内兄名政，字存周，现任工部员外郎，其为人谦恭厚道，
        大有祖父遗风，非膏粱轻薄之流。}
}
{
    \eA{be able to give effect to my modest wishes. As for any outlay that may prove
        necessary, I have given proper explanation, in the letter to my
        brother-in-law, so that you, my brother, need not trouble yourself by giving
        way to much anxiety." As Yü-ts'un bowed and expressed his appreciation in
        most profuse language,-- "Pray,"}
}
{
}

\Verse{12}
{
    \zA{故弟致书烦托，否则不但有污尊兄清操，即弟亦 不屑为矣。”}
}
{
    \eA{he asked, "where does your honoured brother-in-law reside? and what is his
        official capacity? But I fear I'm too coarse in my manner, and could not
        presume to obtrude myself in his presence."}
}
{
}

\Verse{13}
{
    \zA{雨村听了，心下方信了昨日子兴之言，于是又谢了林如海。}
}
{
    \eA{Ju-hai smiled. "And yet," he remarked, "this brother-in-law of mine is after
        all of one and the same family as your worthy self, for he is the grandson
        of the Duke Jung. My elder brother-in-law has now inherited the status of
        Captain-General of the first grade.}
}
{
}

\Verse{14}
{
    \zA{如海又说： “择了出月初二日小女入都，吾兄即同路而往，岂不两便？”}
}
{
    \eA{His name is She, his style Ngen-hou. My second brother-in-law's name is
        Cheng, his style is Tzu-chou. His present post is that of a Second class
        Secretary in the Board of Works. He is modest and kindhearted,}
}
{
}

\Verse{15}
{
    \zA{雨村唯唯听命，心中 十分得意。}
}
{
    \eA{and has much in him of the habits of his grandfather; not one of that
        purse-proud and haughty kind of men.}
}
{
}

\Verse{16}
{
    \zA{如海遂打点礼物并饯行之事，雨村一一领了。}
}
{
    \eA{That is why I have written to him and made the request on your behalf. Were
        he different to what he really is, not only would he cast a slur upon your
        honest purpose,}
}
{
}

\Verse{17}
{
    \zA{那女学生原不忍离亲而去，无奈他外祖母必欲其往，且兼如海说：“汝父年已 半百，再无续室之意，且汝多病，年又极小，上无亲母教养，下无姊妹扶持。}
}
{
    \eA{honourable brother, but I myself likewise would not have been as prompt in
        taking action." When Yü-ts'un heard these remarks, he at length credited
        what had been told him by Tzu-hsing the day before, and he lost no time in
        again expressing his sense of gratitude to Lin Ju-hai. Ju-hai resumed the
        conversation. "I have fixed," (he explained,)}
}
{
}

\Verse{18}
{
    \zA{今去 依傍外祖母及舅氏姊妹，正好减我内顾之忧，如何不去？”}
}
{
    \eA{"upon the second of next month, for my young daughter's departure for the
        capital, and, if you, brother mine, were to travel along with her, would it
        not be an advantage to herself,}
}
{
}

\Verse{19}
{
    \zA{黛玉听了，方洒泪拜别， 随了奶娘及荣府中几个老妇登舟而去。}
}
{
    \eA{as well as to yourself?" Yü-ts'un signified his acquiescence as he listened
        to his proposal; feeling in his inner self extremely elated. Ju-hai availed
        himself of the earliest opportunity to get ready the presents (for the
        capital) and all the requirements for the journey,}
}
{
}

\Verse{20}
{
    \zA{雨村另有船只，带了两个小童，依附黛玉而 行。}
}
{
    \eA{which (when completed,) Yü-ts'un took over one by one. His pupil could not,
        at first, brook the idea, of a separation from her father,}
}
{
}

\Verse{21}
{
    \zA{一日到了京都，雨村先整了衣冠，带着童仆，拿了宗侄的名帖至荣府门上投了。}
}
{
    \eA{but the pressing wishes of her grandmother left her no course (but to
        comply). "Your father," Ju-hai furthermore argued with her, "is already
        fifty; and I entertain no wish to marry again; and then you are always
        ailing;}
}
{
}

\Verse{22}
{
    \zA{彼时贾政已看了妹丈之书，即忙请入相会。}
}
{
    \eA{besides, with your extreme youth, you have, above, no mother of your own to
        take care of you, and below, no sisters to attend to you.}
}
{
}

\Verse{23}
{
    \zA{见雨村像貌魁伟，言谈不俗，且这贾政 最喜的是读书人，礼贤下士，拯溺救危，大有祖风，况又系妹丈致意，因此优待雨 村，更又不同。}
}
{
    \eA{If you now go and have your maternal grandmother, as well as your mother's
        brothers and your cousins to depend upon, you will be doing the best thing
        to reduce the anxiety which I feel in my heart on your behalf. Why then
        should you not go?" Tai-yü, after listening to what her father had to say,
        parted from him in a flood of tears and followed her nurse and several old
        matrons from the Jung mansion on board her boat,}
}
{
}

\Verse{24}
{
    \zA{便极力帮助，题奏之日，谋了一个复职。}
}
{
    \eA{and set out on her journey. Yü-ts'un had a boat to himself, and with two
        youths to wait on him, he prosecuted his voyage in the wake of Tai-yü.}
}
{
}

\Verse{25}
{
    \zA{不上两月，便选了金陵应 天府，辞了贾政，择日到任去了，不在话下。}
}
{
    \eA{By a certain day, they reached Ching Tu; and Yü-ts'un, after first adjusting
        his hat and clothes, came, attended by a youth, to the door of the Jung
        mansion, and sent in a card, which showed his lineage. Chia Cheng had, by
        this time,}
}
{
}

\Verse{26}
{
    \zA{且说黛玉自那日弃舟登岸时，便有荣府打发轿子并拉行李车辆伺候。}
}
{
    \eA{perused his brother-in-law's letter, and he speedily asked him to walk in.
        When they met, he found in Yü-ts'un an imposing manner and polite address.
        This Chia Cheng had, in fact, a great penchant above all things for men of
        education,}
}
{
}

\Verse{27}
{
    \zA{这黛玉尝 听得母亲说，他外祖母家与别人家不同。}
}
{
    \eA{men courteous to the talented, respectful to the learned, ready to lend a
        helping hand to the needy and to succour the distressed, and was, to a great
        extent,}
}
{
}

\Verse{28}
{
    \zA{他近日所见的这几个三等的仆妇，吃穿用 度已是不凡，何况今至其家，都要步步留心，时时在意，不要多说一句话，不可多 行一步路，恐被人耻笑了去。}
}
{
    \eA{like his grandfather. As it was besides a wish intimated by his
        brother-in-law, he therefore treated Yü-ts'un with a consideration still
        more unusual, and readily strained all his resources to assist him. On the
        very day on which the memorial was submitted to the Throne, he obtained by
        his efforts, a reinstatement to office, and before the expiry of two months,}
}
{
}

\Verse{29}
{
    \zA{自上了轿，进了城，从纱窗中瞧了一瞧，其街市之繁 华，人烟之阜盛，自非别处可比。}
}
{
    \eA{Yü-t'sun was forthwith selected to fill the appointment of prefect of Ying
        T'ien in Chin Ling. Taking leave of Chia Cheng, he chose a propitious day,
        and proceeded to his post, where we will leave him without further notice
        for the present. But to return to Tai-yü.}
}
{
}

\Verse{30}
{
    \zA{又行了半日，忽见街北蹲着两个大石狮子，三间 兽头大门，门前列坐着十来个华冠丽服之人，正门不开，只东西两角门有人出入。}
}
{
    \eA{On the day on which she left the boat, and the moment she put her foot on
        shore, there were forthwith at her disposal chairs for her own use, and
        carts for the luggage, sent over from the Jung mansion. Lin Tai-yü had often
        heard her mother recount how different was her grandmother's house from that
        of other people's; and having seen for herself how above the common run were
        already the attendants of the three grades,}
}
{
}

\Verse{31}
{
    \zA{正门之上有一匾，匾上大书“敕造宁国府”五个大字。}
}
{
    \eA{(sent to wait upon her,) in attire, in their fare, in all their articles of
        use, "how much more," (she thought to herself) "now that I am going to her
        home,}
}
{
}

\Verse{32}
{
    \zA{黛玉想道：“这是外祖的长 房了。”}
}
{
    \eA{must I be careful at every step, and circumspect at every moment! Nor must I
        utter one word too many,}
}
{
}

\Verse{33}
{
    \zA{又往西不远，照样也是三间大门，方是“荣国府”，却不进正门，只由西 角门而进。}
}
{
    \eA{nor make one step more than is proper, for fear lest I should be ridiculed
        by any of them!" From the moment she got into the chair, and they had
        entered within the city walls, she found, as she looked around, through the
        gauze window,}
}
{
}

\Verse{34}
{
    \zA{轿子抬着走了一箭之远，将转弯时便歇了轿，后面的婆子也都下来了， 另换了四个眉目秀洁的十七八岁的小厮上来，抬着轿子，众婆子步下跟随。}
}
{
    \eA{at the bustle in the streets and public places and at the immense concourse
        of people, everything naturally so unlike what she had seen elsewhere. After
        they had also been a considerable time on the way, she suddenly caught
        sight, at the northern end of the street, of two huge squatting lions of
        marble and of three lofty gates with (knockers representing) the heads of
        animals.}
}
{
}

\Verse{35}
{
    \zA{至一垂 花门前落下，那小斯俱肃然退出，众婆子上前打起轿帘，扶黛玉下了轿。}
}
{
    \eA{In front of these gates, sat, in a row, about ten men in coloured hats and
        fine attire. The main gate was not open. It was only through the side gates,
        on the east and west, that people went in and came out. Above the centre
        gate was a tablet.}
}
{
}

\Verse{36}
{
    \zA{黛玉扶着 婆子的手进了垂花门，两边是超手游廊，正中是穿堂，当地放着一个紫檀架子大理 石屏风。}
}
{
    \eA{On this tablet were inscribed in five large characters--"The Ning Kuo
        mansion erected by imperial command." "This must be grandmother's eldest
        son's residence," reflected Tai-yü. Towards the east, again, at no great
        distance, were three more high gateways, likewise of the same kind as those
        she had just seen.}
}
{
}

\Verse{37}
{
    \zA{转过屏风，小小三间厅房，厅后便是正房大院。}
}
{
    \eA{This was the Jung Kuo mansion. They did not however go in by the main gate;
        but simply made their entrance through the east side door. With the sedans
        on their shoulders,}
}
{
}

\Verse{38}
{
    \zA{正面五间上房，皆是雕梁 画栋，两边穿山游廊厢房，挂着各色鹦鹉画眉等雀鸟。}
}
{
    \eA{(the bearers) proceeded about the distance of the throw of an arrow, when
        upon turning a corner, they hastily put down the chairs. The matrons, who
        came behind, one and all also dismounted. (The bearers) were changed for
        four youths of seventeen or eighteen,}
}
{
}

\Verse{39}
{
    \zA{台阶上坐着几个穿红着绿的 丫头，一见他们来了，都笑迎上来道：“刚才老太太还念诵呢!可巧就来了。”}
}
{
    \eA{with hats and clothes without a blemish, and while they carried the chair,
        the whole bevy of matrons followed on foot. When they reached a
        creeper-laden gate, the sedan was put down, and all the youths stepped back
        and retired. The matrons came forward, raised the screen, and supported
        Tai-yü to descend from the chair.}
}
{
}

\Verse{40}
{
    \zA{于 是三四人争着打帘子。}
}
{
    \eA{Lin Tai-yü entered the door with the creepers, resting on the hand of a
        matron.}
}
{
}

\Verse{41}
{
    \zA{一面听得人说：“林姑娘来了！”}
}
{
    \eA{On both sides was a verandah, like two outstretched arms. An Entrance Hall
        stood in the centre,}
}
{
}

\Verse{42}
{
    \zA{黛玉方进房，只见两个人扶着一位鬓发如银的老母迎上来。}
}
{
    \eA{in the middle of which was a door-screen of Ta Li marble, set in an ebony
        frame. On the other side of this screen were three very small halls. At the
        back of these came at once an extensive courtyard, belonging to the main
        building.}
}
{
}

\Verse{43}
{
    \zA{黛玉知是外祖母 了，正欲下拜，早被外祖母抱住，搂入怀中，“心肝儿肉”叫着大哭起来。}
}
{
    \eA{In the front part were five parlours, the frieze of the ceiling of which was
        all carved, and the pillars ornamented. On either side, were covered
        avenues, resembling passages through a rock. In the side-rooms were
        suspended cages,}
}
{
}

\Verse{44}
{
    \zA{当下侍 立之人无不下泪，黛玉也哭个不休。}
}
{
    \eA{full of parrots of every colour, thrushes, and birds of every description.
        On the terrace-steps, sat several waiting maids, dressed in red and green,}
}
{
}

\Verse{45}
{
    \zA{众人慢慢解劝，那黛玉方拜见了外祖母。}
}
{
    \eA{and the whole company of them advanced, with beaming faces, to greet them,
        when they saw the party approach. "Her venerable ladyship,"}
}
{
}

\Verse{46}
{
    \zA{贾母 方一一指与黛玉道：“这是你大舅母。}
}
{
    \eA{they said, "was at this very moment thinking of you, miss, and, by a strange
        coincidence, here you are." Three or four of them forthwith vied with each
        other in raising the door curtain,}
}
{
}

\Verse{47}
{
    \zA{这是二舅母。}
}
{
    \eA{while at the same time was heard some one announce:}
}
{
}

\Verse{48}
{
    \zA{这是你先前珠大哥的媳妇珠大 嫂子。”}
}
{
    \eA{"Miss Lin has arrived." No sooner had she entered the room, than she espied
        two servants supporting a venerable lady,}
}
{
}

\Verse{49}
{
    \zA{黛玉一一拜见。}
}
{
    \eA{with silver-white hair, coming forward to greet her.}
}
{
}

\Verse{50}
{
    \zA{贾母又叫：“请姑娘们。}
}
{
    \eA{Convinced that this lady must be her grandmother, she was about to prostrate
        herself and pay her obeisance,}
}
{
}

\Verse{51}
{
    \zA{今日远客来了，可以不必上学去。”}
}
{
    \eA{when she was quickly clasped in the arms of her grandmother, who held her
        close against her bosom;}
}
{
}

\Verse{52}
{
    \zA{众人答应了一声，便去了两个。}
}
{
    \eA{and as she called her "my liver! my flesh!" (my love! my darling!) she began
        to sob aloud. The bystanders too,}
}
{
}

\Verse{53}
{
    \zA{不一时，只见三个奶妈并五六个丫鬟，拥着三位姑 娘来了。}
}
{
    \eA{at once, without one exception, melted into tears; and Tai-yü herself found
        some difficulty in restraining her sobs. Little by little the whole party
        succeeded in consoling her,}
}
{
}

\Verse{54}
{
    \zA{第一个肌肤微丰，身材合中，腮凝新荔，鼻腻鹅脂，温柔沉默，观之可亲。}
}
{
    \eA{and Tai-yü at length paid her obeisance to her grandmother. Her ladyship
        thereupon pointed them out one by one to Tai-yü. "This," she said, "is the
        wife of your uncle, your mother's elder brother; this is the wife of your
        uncle,}
}
{
}

\Verse{55}
{
    \zA{第二个削肩细腰，长挑身材，鸭蛋脸儿，俊眼修眉，顾盼神飞，文彩精华，见之忘 俗。}
}
{
    \eA{her second brother; and this is your eldest sister-in-law Chu, the wife of
        your senior cousin Chu." Tai-yü bowed to each one of them (with folded
        arms). "Ask the young ladies in," dowager lady Chia went on to say; "tell
        them a guest from afar has just arrived,}
}
{
}

\Verse{56}
{
    \zA{第三个身量未足，形容尚小。}
}
{
    \eA{one who comes for the first time; and that they may not go to their
        lessons." The servants with one voice signified their obedience,}
}
{
}

\Verse{57}
{
    \zA{其钗环裙袄，三人皆是一样的妆束。}
}
{
    \eA{and two of them speedily went to carry out her orders. Not long after three
        nurses and five or six waiting-maids were seen ushering in three young
        ladies.}
}
{
}

\Verse{58}
{
    \zA{黛玉忙起身 迎上来见礼，互相厮认，归了坐位。}
}
{
    \eA{The first was somewhat plump in figure and of medium height; her cheeks had
        a congealed appearance, like a fresh lichee; her nose was glossy like goose
        fat. She was gracious, demure,}
}
{
}
\ChapterImage{img/chapters/011第3回 雨村依附黛玉進京.jpg}


\Verse{59}
{
    \zA{丫鬟送上茶来。}
}
{
    \eA{and lovable to look at. The second had sloping shoulders,}
}
{
}

\Verse{60}
{
    \zA{不过叙些黛玉之母如何得病， 如何请医服药，如何送死发丧。}
}
{
    \eA{and a slim waist. Tall and slender was she in stature, with a face like the
        egg of a goose. Her eyes so beautiful, with their well-curved eyebrows,
        possessed in their gaze a bewitching flash.}
}
{
}

\Verse{61}
{
    \zA{不免贾母又伤感起来，因说：“我这些女孩儿，所 疼的独有你母亲。}
}
{
    \eA{At the very sight of her refined and elegant manners all idea of vulgarity
        was forgotten. The third was below the medium size, and her mien was, as
        yet, childlike. In their head ornaments,}
}
{
}

\Verse{62}
{
    \zA{今一旦先我而亡，不得见面，怎不伤心！”}
}
{
    \eA{jewelry, and dress, the get-up of the three young ladies was identical.
        Tai-yü speedily rose to greet them and to exchange salutations.}
}
{
}

\Verse{63}
{
    \zA{说着携了黛玉的手又 哭起来。}
}
{
    \eA{After they had made each other's acquaintance, they all took a seat,
        whereupon the servants brought the tea.}
}
{
}

\Verse{64}
{
    \zA{众人都忙相劝慰，方略略止住。}
}
{
    \eA{Their conversation was confined to Tai-yü's mother,--how she had fallen ill,
        what doctors had attended her, what medicines had been given her,}
}
{
}

\Verse{65}
{
    \zA{众人见黛玉年纪虽小，其举止言谈不俗，身体面貌虽弱不胜衣，却有一段风流 态度，便知他有不足之症。}
}
{
    \eA{and how she had been buried and mourned; and dowager lady Chia was naturally
        again in great anguish. "Of all my daughters," she remarked, "your mother
        was the one I loved best, and now in a twinkle, she has passed away, before
        me too, and I've not been able to so much as see her face.}
}
{
}

\Verse{66}
{
    \zA{因问：“常服何药？}
}
{
    \eA{How can this not make my heart sore-stricken?" And as she gave vent to these
        feelings,}
}
{
}

\Verse{67}
{
    \zA{为何不治好了？”}
}
{
    \eA{she took Tai-yü's hand in hers, and again gave way to sobs;}
}
{
}

\Verse{68}
{
    \zA{黛玉道：“我自 来如此，从会吃饭时便吃药，到如今了，经过多少名医，总未见效。}
}
{
    \eA{and it was only after the members of the family had quickly made use of much
        exhortation and coaxing, that they succeeded, little by little, in stopping
        her tears. They all perceived that Tai-yü, despite her youthful years and
        appearance,}
}
{
}

\Verse{69}
{
    \zA{那一年我才三 岁，记得来了一个癞头和尚，说要化我去出家。}
}
{
    \eA{was lady-like in her deportment and address, and that though with her
        delicate figure and countenance, (she seemed as if) unable to bear the very
        weight of her clothes, she possessed, however, a certain captivating air.}
}
{
}

\Verse{70}
{
    \zA{我父母自是不从，他又说：‘既舍 不得他，但只怕他的病一生也不能好的!若要好时，除非从此以后总不许见哭声，
        除父母之外，凡有外亲一概不见，方可平安了此一生。’}
}
{
    \eA{And as they readily noticed the symptoms of a weak constitution, they went
        on in consequence to make inquiries as to what medicines she ordinarily
        took, and how it was that her complaint had not been cured. "I have,"
        explained Tai-yü, "been in this state ever since I was born; though I've
        taken medicines from the very time I was able to eat rice, up to the
        present, and have been treated by ever so many doctors of note,}
}
{
}

\Verse{71}
{
    \zA{这和尚疯疯癫癫说了这些 不经之谈，也没人理他。}
}
{
    \eA{I've not derived any benefit. In the year when I was yet only three, I
        remember a mangy-headed bonze coming to our house, and saying that he would
        take me along,}
}
{
}

\Verse{72}
{
    \zA{如今还是吃人参养荣丸。”}
}
{
    \eA{and make a nun of me; but my father and mother would, on no account, give
        their consent.}
}
{
}

\Verse{73}
{
    \zA{贾母道：“这正好，我这里正 配丸药呢，叫他们多配一料就是了。”}
}
{
    \eA{'As you cannot bear to part from her and to give her up,' he then remarked,
        'her ailment will, I fear, never, throughout her life, be cured. If you wish
        to see her all right,}
}
{
}

\Verse{74}
{
    \zA{一语未完，只听后院中有笑语声，说：“我来迟了，没得迎接远客！”}
}
{
    \eA{it is only to be done by not letting her, from this day forward, on any
        account, listen to the sound of weeping, or see, with the exception of her
        parents, any relatives outside the family circle.}
}
{
}

\Verse{75}
{
    \zA{黛玉思 忖道：“这些人个个皆敛声屏气如此，这来者是谁，这样放诞无礼？”}
}
{
    \eA{Then alone will she be able to go through this existence in peace and in
        quiet.' No one heeded the nonsensical talk of this raving priest; but here
        am I, up to this very day, dosing myself with ginseng pills as a tonic."
        "What a lucky coincidence!"}
}
{
}

\Verse{76}
{
    \zA{心下想时， 只见一群媳妇丫鬟拥着一个丽人从后房进来。}
}
{
    \eA{interposed dowager lady Chia; "some of these pills are being compounded
        here, and I'll simply tell them to have an extra supply made; that's all."
        Hardly had she finished these words,}
}
{
}

\Verse{77}
{
    \zA{这个人打扮与姑娘们不同，彩绣辉 煌，恍若神妃仙子。}
}
{
    \eA{when a sound of laughter was heard from the back courtyard. "Here I am too
        late!" the voice said, "and not in time to receive the distant visitor!"
        "Every one of all these people,"}
}
{
}

\Verse{78}
{
    \zA{头上戴着金丝八宝攒珠髻，绾着朝阳五凤挂珠钗，项上戴着赤 金盘螭缨络圈，身上穿着缕金百蝶穿花大红云缎窄？}
}
{
    \eA{reflected Tai-yü, "holds her peace and suppresses the very breath of her
        mouth; and who, I wonder, is this coming in this reckless and rude manner?"
        While, as yet, preoccupied with these thoughts, she caught sight of a crowd
        of married women and waiting-maids enter from the back room, pressing round
        a regular beauty. The attire of this person bore no similarity to that of
        the young ladies.}
}
{
}

\Verse{79}
{
    \zA{袄，外罩五彩刻丝石青银鼠 褂，下着翡翠撒花洋绉裙。}
}
{
    \eA{In all her splendour and lustre, she looked like a fairy or a goddess. In
        her coiffure, she had a band of gold filigree work, representing the eight
        precious things,}
}
{
}

\Verse{80}
{
    \zA{一双丹凤三角眼，两弯柳叶掉梢眉，身量苗条，体格风 骚，粉面含春威不露，丹唇未启笑先闻。}
}
{
    \eA{inlaid with pearls; and wore pins, at the head of each of which were five
        phoenixes in a rampant position, with pendants of pearls. On her neck, she
        had a reddish gold necklet, like coiled dragons, with a fringe of tassels.
        On her person, she wore a tight-sleeved jacket,}
}
{
}

\Verse{81}
{
    \zA{黛玉连忙起身接见。}
}
{
    \eA{of dark red flowered satin, covered with hundreds of butterflies,}
}
{
}

\Verse{82}
{
    \zA{贾母笑道：“你不认 得他：他是我们这里有名的一个泼辣货，南京所谓‘辣子’，你只叫他‘凤辣子’ 就是了。”}
}
{
    \eA{embroidered in gold, interspersed with flowers. Over all, she had a
        variegated stiff-silk pelisse, lined with slate-blue ermine; while her
        nether garments consisted of a jupe of kingfisher-colour foreign crepe,
        brocaded with flowers. She had a pair of eyes, triangular in shape like
        those of the red phoenix,}
}
{
}

\Verse{83}
{
    \zA{黛玉正不知以何称呼，众姊妹都忙告诉黛玉道：“这是琏二嫂子。”}
}
{
    \eA{two eyebrows, curved upwards at each temple, like willow leaves. Her stature
        was elegant; her figure graceful; her powdered face like dawning spring,
        majestic, yet not haughty. Her carnation lips, long before they parted,}
}
{
}

\Verse{84}
{
    \zA{黛 玉虽不曾识面，听见他母亲说过：大舅贾赦之子贾琏，娶的就是二舅母王氏的内侄 女；自幼假充男儿教养，学名叫做王熙凤。}
}
{
    \eA{betrayed a smile. Tai-yü eagerly rose and greeted her. Old lady Chia then
        smiled. "You don't know her," she observed. "This is a cunning vixen, who
        has made quite a name in this establishment! In Nanking, she went by the
        appellation of vixen, and if you simply call her Feng Vixen, it will do."
        Tai-yü was just at a loss how to address her,}
}
{
}

\Verse{85}
{
    \zA{黛玉忙陪笑见礼，以“嫂”呼之。}
}
{
    \eA{when all her cousins informed Tai-yü, that this was her sister-in-law Lien.
        Tai-yü had not, it is true,}
}
{
}

\Verse{86}
{
    \zA{这熙凤携着黛玉的手，上下细细打量一回，便仍送至贾母身边坐下，因笑道： “天下真有这样标致人儿!我今日才算看见了!况且这通身的气派竟不像老祖宗的
        外孙女儿，竟是嫡亲的孙女儿似的，怨不得老祖宗天天嘴里心里放不下。}
}
{
    \eA{made her acquaintance before, but she had heard her mother mention that her
        eldest maternal uncle Chia She's son, Chia Lien, had married the niece of
        Madame Wang, her second brother's wife, a girl who had, from her infancy,
        purposely been nurtured to supply the place of a son, and to whom the school
        name of Wang Hsi-feng had been given.}
}
{
}

\Verse{87}
{
    \zA{只可怜我 这妹妹这么命苦，怎么姑妈偏就去世了呢！”}
}
{
    \eA{Tai-yü lost no time in returning her smile and saluting her with all
        propriety, addressing her as my sister-in-law. This Hsi-feng laid hold of
        Tai-yü's hand,}
}
{
}

\Verse{88}
{
    \zA{说着便用帕拭泪。}
}
{
    \eA{and minutely scrutinised her, for a while, from head to foot;}
}
{
}

\Verse{89}
{
    \zA{贾母笑道：“我才 好了，你又来招我。}
}
{
    \eA{after which she led her back next to dowager lady Chia, where they both took
        a seat. "If really there be a being of such beauty in the world,"}
}
{
}

\Verse{90}
{
    \zA{你妹妹远路才来，身子又弱，也才劝住了，快别再提了。”}
}
{
    \eA{she consequently observed with a smile, "I may well consider as having set
        eyes upon it to-day! Besides, in the air of her whole person, she doesn't in
        fact look like your granddaughter-in-law,}
}
{
}

\Verse{91}
{
    \zA{熙 凤听了，忙转悲为喜道：“正是呢!我一见了妹妹，一心都在他身上，又是喜欢， 又是伤心，竟忘了老祖宗了，该打，该打！”}
}
{
    \eA{our worthy ancestor, but in every way like your ladyship's own kindred-
        granddaughter! It's no wonder then that your venerable ladyship should have,
        day after day, had her unforgotten, even for a second, in your lips and
        heart. It's a pity, however, that this cousin of mine should have such a
        hard lot! How did it happen that our aunt died at such an early period?"}
}
{
}

\Verse{92}
{
    \zA{又忙拉着黛玉的手问道：“妹妹几岁 了？}
}
{
    \eA{As she uttered these words, she hastily took her handkerchief and wiped the
        tears from her eyes. "I've only just recovered from a fit of crying,"
        dowager lady Chia observed,}
}
{
}

\Verse{93}
{
    \zA{可也上过学？}
}
{
    \eA{as she smiled, "and have you again come to start me?}
}
{
}

\Verse{94}
{
    \zA{现吃什么药？}
}
{
    \eA{Your cousin has only now arrived from a distant journey,}
}
{
}

\Verse{95}
{
    \zA{在这里别想家，要什么吃的、什么玩的，只管告诉我。}
}
{
    \eA{and she is so delicate to boot! Besides, we have a few minutes back
        succeeded in coaxing her to restrain her sobs, so drop at once making any
        allusion to your former remarks!"}
}
{
}

\Verse{96}
{
    \zA{丫头老婆们不好，也只管告诉我。”}
}
{
    \eA{This Hsi-feng, upon hearing these words, lost no time in converting her
        sorrow into joy. "Quite right," she remarked.}
}
{
}

\Verse{97}
{
    \zA{黛玉一一答应。}
}
{
    \eA{"But at the sight of my cousin, my whole heart was absorbed in her,}
}
{
}

\Verse{98}
{
    \zA{一面熙凤又问人：“林姑娘的 东西可搬进来了？}
}
{
    \eA{and I felt happy, and yet wounded at heart: but having disregarded my
        venerable ancestor's presence, I deserve to be beaten, I do indeed!" And
        hastily taking once more Tai-yü's hand in her own:}
}
{
}

\Verse{99}
{
    \zA{带了几个人来？}
}
{
    \eA{"How old are you, cousin?" she inquired; "Have you been to school?}
}
{
}

\Verse{100}
{
    \zA{你们赶早打扫两间屋子，叫他们歇歇儿去。”}
}
{
    \eA{What medicines are you taking? while you live here, you mustn't feel
        homesick; and if there's anything you would like to eat, or to play with,
        mind you come and tell me!}
}
{
}

\Verse{101}
{
    \zA{说话 时已摆了果茶上来，熙凤亲自布让。}
}
{
    \eA{or should the waiting maids or the matrons fail in their duties, don't
        forget also to report them to me." Addressing at the same time the matrons,}
}
{
}

\Verse{102}
{
    \zA{又见二舅母问他：“月钱放完了没有？”}
}
{
    \eA{she went on to ask, "Have Miss Lin's luggage and effects been brought in?
        How many servants has she brought along with her? Go, as soon as you can,
        and sweep two lower rooms and ask them to go and rest."}
}
{
}

\Verse{103}
{
    \zA{熙凤 道：“放完了。}
}
{
    \eA{As she spake, tea and refreshments had already been served,}
}
{
}

\Verse{104}
{
    \zA{刚才带了人到后楼上找缎子，找了半日也没见昨儿太太说的那个。}
}
{
    \eA{and Hsi-feng herself handed round the cups and offered the fruits. Upon
        hearing the question further put by her maternal aunt Secunda, "Whether the
        issue of the monthly allowances of money had been finished or not yet?"}
}
{
}

\Verse{105}
{
    \zA{想必太太记错了。”}
}
{
    \eA{Hsi-feng replied: "The issue of the money has also been completed;}
}
{
}

\Verse{106}
{
    \zA{王夫人道：“有没有，什么要紧。”}
}
{
    \eA{but a few moments back, when I went along with several servants to the back
        upper-loft, in search of the satins,}
}
{
}

\Verse{107}
{
    \zA{因又说道：“该随手拿出 两个来给你这妹妹裁衣裳啊。}
}
{
    \eA{we looked for ever so long, but we saw nothing of the kind of satins alluded
        to by you, madame, yesterday; so may it not be that your memory misgives
        you?" "Whether there be any or not,}
}
{
}

\Verse{108}
{
    \zA{等晚上想着再叫人去拿罢。”}
}
{
    \eA{of that special kind, is of no consequence," observed madame Wang. "You
        should take out,"}
}
{
}

\Verse{109}
{
    \zA{熙凤道：“我倒先料着 了。}
}
{
    \eA{she therefore went on to add, "any two pieces which first come under your
        hand,}
}
{
}

\Verse{110}
{
    \zA{知道妹妹这两日必到，我已经预备下了，等太太回去过了目，好送来。”}
}
{
    \eA{for this cousin of yours to make herself dresses with; and in the evening,
        if I don't forget, I'll send some one to fetch them." "I've in fact already
        made every provision," rejoined Hsi-feng; "knowing very well that my cousin
        would be arriving within these two days,}
}
{
}

\Verse{111}
{
    \zA{王夫 人一笑，点头不语。}
}
{
    \eA{I have had everything got ready for her. And when you, madame, go back,}
}
{
}

\Verse{112}
{
    \zA{当下茶果已撤，贾母命两个老嬷嬷带黛玉去见两个舅舅去。}
}
{
    \eA{if you will pass an eye over everything, I shall be able to send them
        round." Madame Wang gave a smile, nodded her head assentingly, but uttered
        not a word by way of reply.}
}
{
}

\Verse{113}
{
    \zA{维时贾赦之妻邢氏 忙起身笑回道：“我带了外甥女儿过去，到底便宜些。”}
}
{
    \eA{The tea and fruit had by this time been cleared, and dowager lady Chia
        directed two old nurses to take Tai-yü to go and see her two maternal
        uncles; whereupon Chia She's wife, madame Hsing, hastily stood up and with a
        smiling face suggested,}
}
{
}

\Verse{114}
{
    \zA{贾母笑道：“正是呢。}
}
{
    \eA{"I'll take my niece over; for it will after all be considerably better if I
        go!"}
}
{
}

\Verse{115}
{
    \zA{你 也去罢，不必过来了。”}
}
{
    \eA{"Quite so!" answered dowager lady Chia, smiling; "you can go home too,}
}
{
}
\ChapterImage{img/chapters/012第3回 接外孫賈母憐孤女.jpg}


\Verse{116}
{
    \zA{那邢夫人答应了，遂带着黛玉和王夫人作辞，大家送至穿 堂。}
}
{
    \eA{and there will be no need for you to come over again!" Madame Hsing
        expressed her assent, and forthwith led Tai-yü to take leave of madame Wang.
        The whole party escorted them as far as the door of the Entrance Hall,}
}
{
}

\Verse{117}
{
    \zA{垂花门前早有众小厮拉过一辆翠幄清油车来，邢夫人携了黛玉坐上，众老婆们 放下车帘，方命小厮们抬起。}
}
{
    \eA{hung with creepers, where several youths had drawn a carriage, painted light
        blue, with a kingfisher-coloured hood. Madame Hsing led Tai-yü by the hand
        and they got up into their seats. The whole company of matrons put the
        curtain down, and then bade the youths raise the carriage; who dragged it
        along, until they came to an open space,}
}
{
}

\Verse{118}
{
    \zA{拉至宽处，驾上驯骡，出了西角门往东，过荣府正门， 入一黑油漆大门内，至仪门前方下了车。}
}
{
    \eA{where they at length put the mules into harness. Going out again by the
        eastern side gate, they proceeded in an easterly direction, passed the main
        entrance of the Jung mansion, and entered a lofty doorway painted black. On
        the arrival in front of the ceremonial gate,}
}
{
}

\Verse{119}
{
    \zA{邢夫人挽着黛玉的手进入院中，黛玉度其 处必是荣府中之花园隔断过来的。}
}
{
    \eA{they at once dismounted from the curricle, and madame Hsing, hand-in-hand
        with Tai-yü, walked into the court. "These grounds," surmised Tai-yü to
        herself, "must have been originally converted from a piece partitioned from
        the garden of the Jung mansion."}
}
{
}

\Verse{120}
{
    \zA{进入三层仪门，果见正房、厢房、游廊，悉皆小 巧别致，不似那边的轩峻壮丽，且院中随处之树木山石皆好。}
}
{
    \eA{Having entered three rows of ceremonial gates they actually caught sight of
        the main structure, with its vestibules and porches, all of which, though on
        a small scale, were full of artistic and unique beauty. They were nothing
        like the lofty, imposing, massive and luxurious style of architecture on the
        other side,}
}
{
}

\Verse{121}
{
    \zA{及进入正室，早有许 多艳妆丽服之姬妾丫鬟迎着。}
}
{
    \eA{yet the avenues and rockeries, in the various places in the court, were all
        in perfect taste. When they reached the interior of the principal pavilion,}
}
{
}

\Verse{122}
{
    \zA{邢夫人让黛玉坐了，一面令人到外书房中请贾赦。}
}
{
    \eA{a large concourse of handmaids and waiting maids, got up in gala dress, were
        already there to greet them. Madame Hsing pressed Tai-yü into a seat,}
}
{
}

\Verse{123}
{
    \zA{一 时回来说：“老爷说：‘连日身上不好，见了姑娘彼此伤心，暂且不忍相见。}
}
{
    \eA{while she bade some one go into the outer library and request Mr. Chia She
        to come over. In a few minutes the servant returned. "Master," she
        explained, "says: 'that he has not felt quite well for several days,}
}
{
}

\Verse{124}
{
    \zA{劝姑 娘不必伤怀想家，跟着老太太和舅母，是和家里一样的。}
}
{
    \eA{that as the meeting with Miss Lin will affect both her as well as himself,
        he does not for the present feel equal to seeing each other, that he advises
        Miss Lin not to feel despondent or homesick;}
}
{
}

\Verse{125}
{
    \zA{姐妹们虽拙，大家一处作 伴，也可以解些烦闷。}
}
{
    \eA{that she ought to feel quite at home with her venerable ladyship, (her
        grandmother,) as well as her maternal aunts; that her cousins are, it is
        true,}
}
{
}

\Verse{126}
{
    \zA{或有委屈之处，只管说，别外道了才是。’”}
}
{
    \eA{blunt, but that if all the young ladies associated together in one place,
        they may also perchance dispel some dulness; that if ever (Miss Lin) has any
        grievance,}
}
{
}

\Verse{127}
{
    \zA{黛玉忙站起身 来，一一答应了。}
}
{
    \eA{she should at once speak out, and on no account feel a stranger; and
        everything will then be right."}
}
{
}

\Verse{128}
{
    \zA{再坐一刻便告辞，邢夫人苦留吃过饭去。}
}
{
    \eA{Tai-yü lost no time in respectfully standing up, resuming her seat after she
        had listened to every sentence of the message to her. After a while, she
        said goodbye,}
}
{
}

\Verse{129}
{
    \zA{黛玉笑回道：“舅母爱 惜赐饭，原不应辞，只是还要过去拜见二舅舅，恐去迟了不恭，异日再领：望舅母 容谅。”}
}
{
    \eA{and though madame Hsing used every argument to induce her to stay for the
        repast and then leave, Tai-yü smiled and said, "I shouldn't under ordinary
        circumstances refuse the invitation to dinner, which you, aunt, in your love
        kindly extend to me, but I have still to cross over and pay my respects to
        my maternal uncle Secundus; if I went too late, it would,}
}
{
}

\Verse{130}
{
    \zA{邢夫人道：“这也罢了。”}
}
{
    \eA{I fear, be a lack of respect on my part; but I shall accept on another
        occasion.}
}
{
}

\Verse{131}
{
    \zA{遂命两个嬷嬷用方才坐来的车送过去。}
}
{
    \eA{I hope therefore that you will, dear aunt, kindly excuse me." "If such be
        the case," madame Hsing replied, "it's all right." And presently directing
        two nurses to take her niece over,}
}
{
}

\Verse{132}
{
    \zA{于是黛 玉告辞。}
}
{
    \eA{in the carriage, in which they had come a while back,}
}
{
}

\Verse{133}
{
    \zA{邢夫人送至仪门前，又嘱咐了众人几句，眼看着车去了方回来。}
}
{
    \eA{Tai-yü thereupon took her leave; madame Hsing escorting her as far as the
        ceremonial gate, where she gave some further directions to all the company
        of servants. She followed the curricle with her eyes so long as it remained
        in sight, and at length retraced her footsteps.}
}
{
}

\Verse{134}
{
    \zA{一时黛玉进入荣府，下了车，只见一条大甬路直接出大门来。}
}
{
    \eA{Tai-yü shortly entered the Jung Mansion, descended from the carriage, and
        preceded by all the nurses, she at once proceeded towards the east, turned a
        corner, passed through an Entrance Hall,}
}
{
}

\Verse{135}
{
    \zA{众嬷嬷引着便往 东转弯，走过一座东西穿堂、向南大厅之后，仪门内大院落，上面五间大正房，两 边厢房鹿顶，耳门钻山，四通八达，轩昂壮丽，比各处不同。}
}
{
    \eA{running east and west, and walked in a southern direction, at the back of
        the Large Hall. On the inner side of a ceremonial gate, and at the upper end
        of a spacious court, stood a large main building, with five apartments,
        flanked on both sides by out-houses (stretching out) like the antlers on the
        head of deer; side-gates, resembling passages through a hill,}
}
{
}

\Verse{136}
{
    \zA{黛玉便知这方是正内 室。}
}
{
    \eA{establishing a thorough communication all round; (a main building) lofty,
        majestic, solid and grand,}
}
{
}

\Verse{137}
{
    \zA{进入堂屋，抬头迎面先见一个赤金九龙青地大匾，匾上写着斗大三个字，是“荣 禧堂”；后有一行小字：“某年月日书赐荣国公贾源”，又有“万几宸翰”之宝。}
}
{
    \eA{and unlike those in the compound of dowager lady Chia. Tai-yü readily
        concluded that this at last was the main inner suite of apartments. A raised
        broad road led in a straight line to the large gate. Upon entering the Hall,
        and raising her head, she first of all perceived before her a large tablet
        with blue ground, upon which figured nine dragons of reddish gold.}
}
{
}

\Verse{138}
{
    \zA{大紫檀雕螭案上设着三尺多高青绿古铜鼎，悬着待漏随朝墨龙大画，一边是錾金 彝，一边是玻璃盆。}
}
{
    \eA{The inscription on this tablet consisted of three characters as large as a
        peck-measure, and declared that this was the Hall of Glorious Felicity. At
        the end, was a row of characters of minute size, denoting the year, month
        and day, upon which His Majesty had been pleased to confer the tablet upon
        Chia Yuan,}
}
{
}

\Verse{139}
{
    \zA{地下两溜十六张楠木圈椅。}
}
{
    \eA{Duke of Jung Kuo. Besides this tablet, were numberless costly articles
        bearing the autograph of the Emperor.}
}
{
}

\Verse{140}
{
    \zA{又有一副对联，乃是乌木联牌镶着 錾金字迹，道是： 座上珠玑昭日月， 堂前黼黻焕烟霞。}
}
{
    \eA{On the large black ebony table, engraved with dragons, were placed three
        antique blue and green bronze tripods, about three feet in height. On the
        wall hung a large picture representing black dragons, such as were seen in
        waiting chambers of the Sui dynasty.}
}
{
}

\Verse{141}
{
    \zA{下面一行小字是：“世教弟勋袭东安郡王穆莳拜手书。”}
}
{
    \eA{On one side stood a gold cup of chased work, while on the other, a crystal
        casket. On the ground were placed, in two rows, sixteen chairs, made of
        hard-grained cedar.}
}
{
}

\Verse{142}
{
    \zA{原来王夫人时常居坐宴息 也不在这正室中，只在东边的三间耳房内。}
}
{
    \eA{There was also a pair of scrolls consisting of black-wood antithetical
        tablets, inlaid with the strokes of words in chased gold. Their burden was
        this: On the platform shine resplendent pearls like sun or moon, And the
        sheen of the Hall façade gleams like russet sky.}
}
{
}

\Verse{143}
{
    \zA{于是嬷嬷们引黛玉进东房门来。}
}
{
    \eA{Below, was a row of small characters, denoting that the scroll had been
        written by the hand of Mu Shih,}
}
{
}

\Verse{144}
{
    \zA{临窗大 炕上铺着猩红洋毯，正面设着大红金钱蟒引枕，秋香色金钱蟒大条褥，两边设一对
        梅花式洋漆小几，左边几上摆着文王鼎，鼎旁匙箸香盒，右边几上摆着汝窑美人觚， 里面插着时鲜花草。}
}
{
    \eA{a fellow-countryman and old friend of the family, who, for his meritorious
        services, had the hereditary title of Prince of Tung Ngan conferred upon
        him. The fact is that madame Wang was also not in the habit of sitting and
        resting, in this main apartment, but in three side-rooms on the east, so
        that the nurses at once led Tai-yü through the door of the eastern wing. On
        a stove-couch, near the window,}
}
{
}

\Verse{145}
{
    \zA{地下面西一溜四张大椅，都搭着银红撒花椅搭，底下四副脚踏； 两边又有一对高几，几上茗碗瓶花俱备。}
}
{
    \eA{was spread a foreign red carpet. On the side of honour, were laid deep red
        reclining-cushions, with dragons, with gold cash (for scales), and an oblong
        brown-coloured sitting-cushion with gold-cash-spotted dragons. On the two
        sides, stood one of a pair of small teapoys of foreign lacquer of
        peach-blossom pattern.}
}
{
}

\Verse{146}
{
    \zA{其馀陈设，不必细说。}
}
{
    \eA{On the teapoy on the left, were spread out Wen Wang tripods, spoons,
        chopsticks and scent-bottles.}
}
{
}

\Verse{147}
{
    \zA{老嬷嬷让黛玉上炕 坐。}
}
{
    \eA{On the teapoy on the right, were vases from the Ju Kiln, painted with girls
        of great beauty,}
}
{
}

\Verse{148}
{
    \zA{炕沿上却也有两个锦褥对设。}
}
{
    \eA{in which were placed seasonable flowers; (on it were) also teacups, a tea
        service and the like articles.}
}
{
}

\Verse{149}
{
    \zA{黛玉度其位次，便不上炕，只就东边椅上坐了。}
}
{
    \eA{On the floor on the west side of the room, were four chairs in a row, all of
        which were covered with antimacassars, embroidered with silverish-red
        flowers,}
}
{
}

\Verse{150}
{
    \zA{本房的丫鬟忙捧上茶来。}
}
{
    \eA{while below, at the feet of these chairs, stood four footstools. On either
        side,}
}
{
}

\Verse{151}
{
    \zA{黛玉一面吃了，打量这些丫鬟们妆饰衣裙、举止行动，果 与别家不同。}
}
{
    \eA{was also one of a pair of high teapoys, and these teapoys were covered with
        teacups and flower vases. The other nick-nacks need not be minutely
        described. The old nurses pressed Tai-yü to sit down on the stove-couch;}
}
{
}

\Verse{152}
{
    \zA{茶未吃了，只见一个穿红绫袄青绸掐牙背心的一个丫鬟走来笑道：“太太说： 请林姑娘到那边坐罢。”}
}
{
    \eA{but, on perceiving near the edge of the couch two embroidered cushions,
        placed one opposite the other, she thought of the gradation of seats, and
        did not therefore place herself on the couch, but on a chair on the eastern
        side of the room; whereupon the waiting maids,}
}
{
}

\Verse{153}
{
    \zA{老嬷嬷听了，于是又引黛玉出来，到了东南三间小正房内。}
}
{
    \eA{in attendance in these quarters, hastened to serve the tea. While Tai-yü was
        sipping her tea, she observed the headgear, dress, deportment and manners of
        the several waiting maids,}
}
{
}

\Verse{154}
{
    \zA{正面炕上横设一张炕桌，上面堆着书籍茶具，靠东壁面西设着半旧的青缎靠背引 枕。}
}
{
    \eA{which she really found so unlike what she had seen in other households. She
        had hardly finished her tea, when she noticed a waiting maid approach,
        dressed in a red satin jacket, and a waistcoat of blue satin with scollops.}
}
{
}

\Verse{155}
{
    \zA{王夫人却坐在西边下首，亦是半旧青缎靠背坐褥，见黛玉来了，便往东让。}
}
{
    \eA{"My lady requests Miss Lin to come over and sit with her," she remarked as
        she put on a smile. The old nurses, upon hearing this message, speedily
        ushered Tai-yü again out of this apartment, into the three-roomed small main
        building by the eastern porch.}
}
{
}

\Verse{156}
{
    \zA{黛 玉心中料定这是贾政之位，因见挨炕一溜三张椅子上也搭着半旧的弹花椅袱，黛玉 便向椅上坐了。}
}
{
    \eA{On the stove-couch, situated at the principal part of the room, was placed,
        in a transverse position, a low couch-table, at the upper end of which were
        laid out, in a heap, books and a tea service. Against the partition-wall, on
        the east side, facing the west, was a reclining pillow, made of blue satin,}
}
{
}

\Verse{157}
{
    \zA{王夫人再三让他上炕，他方挨王夫人坐下。}
}
{
    \eA{neither old nor new. Madame Wang, however, occupied the lower seat, on the
        west side, on which was likewise placed a rather shabby blue satin
        sitting-rug,}
}
{
}

\Verse{158}
{
    \zA{王夫人因说：“你舅舅 今日斋戒去了，再见罢。}
}
{
    \eA{with a back-cushion; and upon perceiving Tai-yü come in she urged her at
        once to sit on the east side. Tai-yü concluded, in her mind, that this seat
        must certainly belong to Chia Cheng, and espying, next to the couch,}
}
{
}

\Verse{159}
{
    \zA{只是有句话嘱咐你：你三个姐妹倒都极好，以后一处念书 认字，学针线，或偶一玩笑，却都有个尽让的。}
}
{
    \eA{a row of three chairs, covered with antimacassars, strewn with embroidered
        flowers, somewhat also the worse for use, Tai-yü sat down on one of these
        chairs. But as madame Wang pressed her again and again to sit on the couch,
        Tai-yü had at length to take a seat next to her. "Your uncle," madame Wang
        explained,}
}
{
}

\Verse{160}
{
    \zA{我就只一件不放心：我有一个孽根 祸胎，是家里的‘混世魔王’，今日因往庙里还愿去，尚未回来，晚上你看见就知 道了。}
}
{
    \eA{"is gone to observe this day as a fast day, but you'll see him by and bye.
        There's, however, one thing I want to talk to you about. Your three female
        cousins are all, it is true, everything that is nice; and you will, when
        later on you come together for study, or to learn how to do needlework,}
}
{
}

\Verse{161}
{
    \zA{你以后总不用理会他，你这些姐姐妹妹都不敢沾惹他的。”}
}
{
    \eA{or whenever, at any time, you romp and laugh together, find them all most
        obliging; but there's one thing that causes me very much concern. I have
        here one, who is the very root of retribution,}
}
{
}

\Verse{162}
{
    \zA{黛玉素闻母亲说 过，有个内侄乃衔玉而生，顽劣异常，不喜读书，最喜在内帏厮混，外祖母又溺爱， 无人敢管。}
}
{
    \eA{the incarnation of all mischief, one who is a ne'er-do-well, a prince of
        malignant spirits in this family. He is gone to-day to pay his vows in the
        temple, and is not back yet, but you will see him in the evening, when you
        will readily be able to judge for yourself. One thing you must do, and that
        is,}
}
{
}

\Verse{163}
{
    \zA{今见王夫人所说，便知是这位表兄，一面陪笑道：“舅母所说，可是衔 玉而生的？}
}
{
    \eA{from this time forth, not to pay any notice to him. All these cousins of
        yours don't venture to bring any taint upon themselves by provoking him."
        Tai-yü had in days gone by heard her mother explain that she had a nephew,
        born into the world, holding a piece of jade in his mouth,}
}
{
}

\Verse{164}
{
    \zA{在家时记得母亲常说，这位哥哥比我大一岁，小名就叫宝玉，性虽憨顽， 说待姊妹们却是极好的。}
}
{
    \eA{who was perverse beyond measure, who took no pleasure in his books, and
        whose sole great delight was to play the giddy dog in the inner apartments;
        that her maternal grandmother, on the other hand, loved him so fondly that
        no one ever presumed to call him to account, so that when, in this instance,}
}
{
}

\Verse{165}
{
    \zA{况我来了，自然和姊妹们一处，弟兄们是另院别房，岂有 沾惹之理？”}
}
{
    \eA{she heard madame Wang's advice, she at once felt certain that it must be
        this very cousin. "Isn't it to the cousin born with jade in his mouth, that
        you are alluding to, aunt?" she inquired as she returned her smile.}
}
{
}

\Verse{166}
{
    \zA{王夫人笑道：“你不知道原故：他和别人不同，自幼因老太太疼爱， 原系和姐妹们一处娇养惯了的。}
}
{
    \eA{"When I was at home, I remember my mother telling me more than once of this
        very cousin, who (she said) was a year older than I, and whose infant name
        was Pao-yü. She added that his disposition was really wayward, but that he
        treats all his cousins with the utmost consideration. Besides, now that I
        have come here,}
}
{
}

\Verse{167}
{
    \zA{若姐妹们不理他，他倒还安静些；若一日姐妹们和 他多说了一句话，他心上一喜，便生出许多事来。}
}
{
    \eA{I shall, of course, be always together with my female cousins, while the
        boys will have their own court, and separate quarters; and how ever will
        there be any cause of bringing any slur upon myself by provoking him?" "You
        don't know the reasons (that prompt me to warn you)," replied madame Wang
        laughingly.}
}
{
}

\Verse{168}
{
    \zA{所以嘱咐你别理会他。}
}
{
    \eA{"He is so unlike all the rest, all because he has, since his youth up,}
}
{
}

\Verse{169}
{
    \zA{他嘴里一 时甜言蜜语，一时有天没日，疯疯傻傻，只休信他。”}
}
{
    \eA{been doated upon by our old lady! The fact is that he has been spoilt,
        through over-indulgence, by being always in the company of his female
        cousins! If his female cousins pay no heed to him,}
}
{
}

\Verse{170}
{
    \zA{黛玉一一的都答应着。}
}
{
    \eA{he is, at any rate, somewhat orderly, but the day his cousins say one word
        more to him than usual,}
}
{
}

\Verse{171}
{
    \zA{忽见一个丫鬟来说：“老太太那里传晚饭了。”}
}
{
    \eA{much trouble forthwith arises, at the outburst of delight in his heart.
        That's why I enjoin upon you not to heed him. From his mouth, at one time,}
}
{
}

\Verse{172}
{
    \zA{王夫人忙携了黛玉出后房门， 由后廊往西。}
}
{
    \eA{issue sugared words and mellifluous phrases; and at another, like the
        heavens devoid of the sun, he becomes a raving fool; so whatever you do,
        don't believe all he says."}
}
{
}

\Verse{173}
{
    \zA{出了角门，是一条南北甬路，南边是倒座三间小小抱厦厅，北边立着 一个粉油大影壁，后有一个半大门，小小一所房屋。}
}
{
    \eA{Tai-yü was assenting to every bit of advice as it was uttered, when
        unexpectedly she beheld a waiting-maid walk in. "Her venerable ladyship over
        there," she said, "has sent word about the evening meal." Madame Wang
        hastily took Tai-yü by the hand, and emerging by the door of the back-room,
        they went eastwards by the verandah at the back.}
}
{
}
\ChapterImage{img/chapters/013第3回 林黛玉初至寧榮國府.jpg}


\Verse{174}
{
    \zA{王夫人笑指向黛玉道：“这是 你凤姐姐的屋子。}
}
{
    \eA{Past the side gate, was a roadway, running north and south. On the southern
        side were a pavilion with three divisions and a Reception Hall with a
        colonnade.}
}
{
}

\Verse{175}
{
    \zA{回来你好往这里找他去，少什么东西只管和他说就是了。”}
}
{
    \eA{On the north, stood a large screen wall, painted white; behind it was a very
        small building, with a door of half the ordinary size. "These are your
        cousin Feng's rooms,"}
}
{
}

\Verse{176}
{
    \zA{这院 门上也有几个才总角的小厮，都垂手侍立。}
}
{
    \eA{explained madame Wang to Tai-yü, as she pointed to them smiling. "You'll
        know in future your way to come and find her; and if you ever lack anything,}
}
{
}

\Verse{177}
{
    \zA{王夫人遂携黛玉穿过一个东西穿堂，便 是贾母的后院了。}
}
{
    \eA{mind you mention it to her, and she'll make it all right." At the door of
        this court, were also several youths, who had recently had the tufts of
        their hair tied together,}
}
{
}

\Verse{178}
{
    \zA{于是进入后房门，已有许多人在此伺候，见王夫人来，方安设桌 椅。}
}
{
    \eA{who all dropped their hands against their sides, and stood in a respectful
        posture. Madame Wang then led Tai-yü by the hand through a corridor, running
        east and west, into what was dowager lady Chia's back-court.}
}
{
}

\Verse{179}
{
    \zA{贾珠之妻李氏捧杯，熙凤安箸，王夫人进羹。}
}
{
    \eA{Forthwith they entered the door of the back suite of rooms, where stood,
        already in attendance, a large number of servants,}
}
{
}

\Verse{180}
{
    \zA{贾母正面榻上独坐，两旁四张空 椅。}
}
{
    \eA{who, when they saw madame Wang arrive, set to work setting the tables and
        chairs in order. Chia Chu's wife, née Li,}
}
{
}

\Verse{181}
{
    \zA{熙凤忙拉黛玉在左边第一张椅子上坐下，黛玉十分推让。}
}
{
    \eA{served the eatables, while Hsi-feng placed the chopsticks, and madame Wang
        brought the soup in. Dowager lady Chia was seated all alone on the divan, in
        the main part of the apartment,}
}
{
}

\Verse{182}
{
    \zA{贾母笑道：“你舅母 和嫂子们是不在这里吃饭的。}
}
{
    \eA{on the two sides of which stood four vacant chairs. Hsi-feng at once drew
        Tai-yü, meaning to make her sit in the foremost chair on the left side, but
        Tai-yü steadily and concedingly declined.}
}
{
}

\Verse{183}
{
    \zA{你是客，原该这么坐。”}
}
{
    \eA{"Your aunts and sisters-in-law, standing on the right and left,"}
}
{
}

\Verse{184}
{
    \zA{黛玉方告了坐，就坐了。}
}
{
    \eA{dowager lady Chia smilingly explained, "won't have their repast in here,}
}
{
}

\Verse{185}
{
    \zA{贾 母命王夫人也坐了。}
}
{
    \eA{and as you're a guest, it's but proper that you should take that seat."}
}
{
}

\Verse{186}
{
    \zA{迎春姊妹三个告了坐方上来，迎春坐右手第一，探春左第二， 惜春右第二。}
}
{
    \eA{Then alone it was that Tai-yü asked for permission to sit down, seating
        herself on the chair. Madame Wang likewise took a seat at old lady Chia's
        instance; and the three cousins, Ying Ch'un and the others, having craved
        for leave to sit down,}
}
{
}

\Verse{187}
{
    \zA{旁边丫鬟执着拂尘、漱盂、巾帕，李纨、凤姐立于案边布让；外间伺 候的媳妇丫鬟虽多，却连一声咳嗽不闻。}
}
{
    \eA{at length came forward, and Ying Ch'un took the first chair on the right,
        T'an Ch'un the second, and Hsi Ch'un the second on the left. Waiting maids
        stood by holding in their hands, flips and finger-bowls and napkins, while
        Mrs. Li and lady Feng, the two of them, kept near the table advising them
        what to eat,}
}
{
}

\Verse{188}
{
    \zA{饭毕，各各有丫鬟用小茶盘捧上茶来。}
}
{
    \eA{and pressing them to help themselves. In the outer apartments, the married
        women and waiting-maids in attendance, were, it is true, very numerous;}
}
{
}

\Verse{189}
{
    \zA{当 日林家教女以惜福养身，每饭后必过片时方吃茶，不伤脾胃；今黛玉见了这里许多 规矩，不似家中，也只得随和些，接了茶。}
}
{
    \eA{but not even so much as the sound of the cawing of a crow could be heard.
        The repast over, each one was presented by a waiting-maid, with tea in a
        small tea tray; but the Lin family had all along impressed upon the mind of
        their daughter that in order to show due regard to happiness, and to
        preserve good health, it was essential, after every meal, to wait a while,
        before drinking any tea,}
}
{
}

\Verse{190}
{
    \zA{又有人捧过漱盂来，黛玉也漱了口，又 盥手毕。}
}
{
    \eA{so that it should not do any harm to the intestines. When, therefore, Tai-yü
        perceived how many habits there were in this establishment unlike those
        which prevailed in her home, she too had no alternative but to conform
        herself to a certain extent with them. Upon taking over the cup of tea,
        servants came once more and presented finger-bowls for them to rinse their
        mouths,}
}
{
}

\Verse{191}
{
    \zA{然后又捧上茶来，这方是吃的茶。}
}
{
    \eA{and Tai-yü also rinsed hers; and after they had all again finished washing
        their hands, tea was eventually served a second time,}
}
{
}

\Verse{192}
{
    \zA{贾母便说：“你们去罢，让我们自在说 说话儿。”}
}
{
    \eA{and this was, at length, the tea that was intended to be drunk. "You can all
        go," observed dowager lady Chia, "and let us alone to have a chat."}
}
{
}

\Verse{193}
{
    \zA{王夫人遂起身，又说了两句闲话儿，方引李、凤二人去了。}
}
{
    \eA{Madame Wang rose as soon as she heard these words, and having made a few
        irrelevant remarks, she led the way and left the room along with the two
        ladies, Mrs. Li and lady Feng. Dowager lady Chia,}
}
{
}

\Verse{194}
{
    \zA{贾母因问黛玉念何书。}
}
{
    \eA{having inquired of Tai-yü what books she was reading, "I have just begun
        reading the Four Books,"}
}
{
}

\Verse{195}
{
    \zA{黛玉道：“刚念了《四书》。”}
}
{
    \eA{Tai-yü replied. "What books are my cousins reading?" Tai-yü went on to ask.
        "Books,}
}
{
}

\Verse{196}
{
    \zA{黛玉又问姊妹们读何书， 贾母道：“读什么书，不过认几个字罢了。”}
}
{
    \eA{you say!" exclaimed dowager lady Chia; "why all they know are a few
        characters, that's all." The sentence was barely out of her lips, when a
        continuous sounding of footsteps was heard outside,}
}
{
}

\Verse{197}
{
    \zA{一语未了，只听外面一阵脚步响，丫 鬟进来报道：“宝玉来了。”}
}
{
    \eA{and a waiting maid entered and announced that Pao-yü was coming. Tai-yü was
        speculating in her mind how it was that this Pao-yü had turned out such a
        good-for-nothing fellow, when he happened to walk in.}
}
{
}

\Verse{198}
{
    \zA{黛玉心想，这个宝玉不知是怎样个惫懒人呢。}
}
{
    \eA{He was, in fact, a young man of tender years, wearing on his head, to hold
        his hair together, a cap of gold of purplish tinge, inlaid with precious
        gems.}
}
{
}

\Verse{199}
{
    \zA{及至进 来一看，却是位青年公子：头上戴着束发嵌宝紫金冠，齐眉勒着二龙戏珠金抹额，
        一件二色金百蝶穿花大红箭袖，束着五彩丝攒花结长穗宫绦，外罩石青起花八团倭 缎排穗褂，登着青缎粉底小朝靴。}
}
{
    \eA{Parallel with his eyebrows was attached a circlet, embroidered with gold,
        and representing two dragons snatching a pearl. He wore an archery-sleeved
        deep red jacket, with hundreds of butterflies worked in gold of two
        different shades, interspersed with flowers; and was girded with a sash of
        variegated silk, with clusters of designs, to which was attached long
        tassels;}
}
{
}

\Verse{200}
{
    \zA{面若中秋之月，色如春晓之花，鬓若刀裁，眉如 墨画，鼻如悬胆，睛若秋波，虽怒时而似笑，即？}
}
{
    \eA{a kind of sash worn in the palace. Over all, he had a slate-blue fringed
        coat of Japanese brocaded satin, with eight bunches of flowers in relief;
        and wore a pair of light blue satin white-soled, half-dress court-shoes. His
        face was like the full moon at mid-autumn; his complexion, like morning
        flowers in spring;}
}
{
}

\Verse{201}
{
    \zA{视而有情。}
}
{
    \eA{the hair along his temples,}
}
{
}

\Verse{202}
{
    \zA{项上金螭缨络，又有 一根五色丝绦，系着一块美玉。}
}
{
    \eA{as if chiselled with a knife; his eyebrows, as if pencilled with ink; his
        nose like a suspended gallbladder (a well-cut and shapely nose); his eyes
        like vernal waves;}
}
{
}

\Verse{203}
{
    \zA{黛玉一见便吃一大惊，心中想道：“好生奇怪，倒 像在那里见过的，何等眼熟！”}
}
{
    \eA{his angry look even resembled a smile; his glance, even when stern, was full
        of sentiment. Round his neck he had a gold dragon necklet with a fringe;
        also a cord of variegated silk, to which was attached a piece of beautiful
        jade.}
}
{
}

\Verse{204}
{
    \zA{只见这宝玉向贾母请了安，贾母便命：“去见你娘 来。”}
}
{
    \eA{As soon as Tai-yü became conscious of his presence, she was quite taken
        aback. "How very strange!" she was reflecting in her mind; "it would seem as
        if I had seen him somewhere or other,}
}
{
}

\Verse{205}
{
    \zA{即转身去了。}
}
{
    \eA{for his face appears extremely familiar to my eyes;"}
}
{
}

\Verse{206}
{
    \zA{一回再来时，已换了冠带，头上周围一转的短发都结成小辫， 红丝结束，共攒至顶中胎发，总编一根大辫，黑亮如漆，从顶至梢，一串四颗大珠， 用金八宝坠脚。}
}
{
    \eA{when she noticed Pao-yü face dowager lady Chia and make his obeisance. "Go
        and see your mother and then come back," remarked her venerable ladyship;
        and at once he turned round and quitted the room. On his return, he had
        already changed his hat and suit. All round his head, he had a fringe of
        short hair,}
}
{
}

\Verse{207}
{
    \zA{身上穿着银红撒花半旧大袄，仍旧带着项圈、宝玉、寄名锁、护身 符等物，下面半露松绿撒花绫裤，锦边弹墨袜，厚底大红鞋。}
}
{
    \eA{plaited into small queues, and bound with red silk. The queues were gathered
        up at the crown, and all the hair, which had been allowed to grow since his
        birth, was plaited into a thick queue, which looked as black and as glossy
        as lacquer. Between the crown of the head and the extremity of the queue,}
}
{
}

\Verse{208}
{
    \zA{越显得面如傅粉，唇 若施脂，转盼多情，语言若笑。}
}
{
    \eA{hung a string of four large pearls, with pendants of gold, representing the
        eight precious things. On his person, he wore a long silvery-red coat,}
}
{
}

\Verse{209}
{
    \zA{天然一段风韵，全在眉梢；平生万种情思，悉堆眼 角。}
}
{
    \eA{more or less old, bestrewn with embroidery of flowers. He had still round
        his neck the necklet, precious gem, amulet of Recorded Name, philacteries,
        and other ornaments.}
}
{
}

\Verse{210}
{
    \zA{看其外貌最是极好，却难知其底细，后人有《西江月》二词，批的极确。}
}
{
    \eA{Below were partly visible a fir-cone coloured brocaded silk pair of
        trousers, socks spotted with black designs, with ornamented edges, and a
        pair of deep red, thick-soled shoes. (Got up as he was now,) his face
        displayed a still whiter appearance,}
}
{
}

\Verse{211}
{
    \zA{词曰： 无故寻愁觅恨，有时似傻如狂。}
}
{
    \eA{as if painted, and his eyes as if they were set off with carnation. As he
        rolled his eyes, they brimmed with love. When he gave utterance to speech,}
}
{
}

\Verse{212}
{
    \zA{纵然生得好皮囊，腹内原 来草莽。}
}
{
    \eA{he seemed to smile. But the chief natural pleasing feature was mainly
        centred in the curve of his eyebrows. The ten thousand and one fond
        sentiments,}
}
{
}

\Verse{213}
{
    \zA{潦倒不通庶务，愚顽怕读文章。}
}
{
    \eA{fostered by him during the whole of his existence, were all amassed in the
        corner of his eyes.}
}
{
}

\Verse{214}
{
    \zA{行为偏僻性乖张，那管世人诽谤。}
}
{
    \eA{His outward appearance may have been pleasing to the highest degree, but yet
        it was no easy matter to fathom what lay beneath it.}
}
{
}

\Verse{215}
{
    \zA{又曰： 富贵不知乐业，贫穷难耐凄凉。}
}
{
    \eA{There are a couple of roundelays, composed by a later poet, (after the
        excellent rhythm of the) Hsi Chiang Yueh,}
}
{
}

\Verse{216}
{
    \zA{可怜辜负好时光，于国于家无望。}
}
{
    \eA{which depict Pao-yü in a most adequate manner. The roundelays run as
        follows: To gloom and passion prone,}
}
{
}

\Verse{217}
{
    \zA{天下无能第 一，古今不肖无双。}
}
{
    \eA{without a rhyme, Inane and madlike was he many a time, His outer self,
        forsooth, fine may have been, But one wild,}
}
{
}

\Verse{218}
{
    \zA{寄言纨？}
}
{
    \eA{howling waste his mind within:}
}
{
}

\Verse{219}
{
    \zA{与膏粱：莫效此儿形状! 却说贾母见他进来，笑道：“外客没见就脱了衣裳了，还不去见你妹妹呢。”}
}
{
    \eA{Addled his brain that nothing he could see; A dunce! to read essays so loth
        to be! Perverse in bearing, in temper wayward; For human censure he had no
        regard. When rich, wealth to enjoy he knew not how; When poor, to poverty he
        could not bow.}
}
{
}

\Verse{220}
{
    \zA{宝玉早已看见了一个袅袅婷婷的女儿，便料定是林姑妈之女，忙来见礼。}
}
{
    \eA{Alas! what utter waste of lustrous grace! To state, to family what a
        disgrace! Of ne'er-do-wells below he was the prime, Unfilial like him none
        up to this time. Ye lads, pampered with sumptuous fare and dress,}
}
{
}

\Verse{221}
{
    \zA{归了坐细 看时，真是与众各别。}
}
{
    \eA{Beware! In this youth's footsteps do not press! But to proceed with our
        story.}
}
{
}

\Verse{222}
{
    \zA{只见： 两弯似蹙非蹙笼烟眉，一双似喜非喜含情目。}
}
{
    \eA{"You have gone and changed your clothes," observed dowager lady Chia,
        "before being introduced to the distant guest. Why don't you yet salute your
        cousin?"}
}
{
}

\Verse{223}
{
    \zA{态生两靥之愁，娇袭一身之病。}
}
{
    \eA{Pao-yü had long ago become aware of the presence of a most beautiful young
        lady, who, he readily concluded,}
}
{
}

\Verse{224}
{
    \zA{泪光点点，娇喘微微。}
}
{
    \eA{must be no other than the daughter of his aunt Lin. He hastened to advance
        up to her,}
}
{
}

\Verse{225}
{
    \zA{闲静似娇花照水，行动如弱柳扶风。}
}
{
    \eA{and make his bow; and after their introduction, he resumed his seat, whence
        he minutely scrutinised her features,}
}
{
}

\Verse{226}
{
    \zA{心较比干多一窍，病如 西子胜三分。}
}
{
    \eA{(which he thought) so unlike those of all other girls. Her two arched
        eyebrows, thick as clustered smoke,}
}
{
}

\Verse{227}
{
    \zA{宝玉看罢，笑道：“这个妹妹我曾见过的。”}
}
{
    \eA{bore a certain not very pronounced frowning wrinkle. She had a pair of eyes,
        which possessed a cheerful, and yet one would say, a sad expression,}
}
{
}

\Verse{228}
{
    \zA{贾母笑道：“又胡说了，你何曾见过？”}
}
{
    \eA{overflowing with sentiment. Her face showed the prints of sorrow stamped on
        her two dimpled cheeks.}
}
{
}

\Verse{229}
{
    \zA{宝玉笑道：“虽没见过，却看着面善，心里倒像是远别重逢的一般。”}
}
{
    \eA{She was beautiful, but her whole frame was the prey of a hereditary disease.
        The tears in her eyes glistened like small specks. Her balmy breath was so
        gentle. She was as demure as a lovely flower reflected in the water.}
}
{
}

\Verse{230}
{
    \zA{贾母笑道： “好，好!这么更相和睦了。”}
}
{
    \eA{Her gait resembled a frail willow, agitated by the wind. Her heart, compared
        with that of Pi Kan, had one more aperture of intelligence;}
}
{
}
\ChapterImage{img/chapters/014第3回 賈寶玉初會林黛玉 寶玉癡狂狠摔那玉.jpg}


\Verse{231}
{
    \zA{宝玉便走向黛玉身边坐下，又细细打量一番，因问：“妹妹可曾读书？”}
}
{
    \eA{while her ailment exceeded (in intensity) by three degrees the ailment of
        Hsi-Tzu. Pao-yü, having concluded his scrutiny of her, put on a smile and
        said, "This cousin I have already seen in days gone by." "There you are
        again with your nonsense,"}
}
{
}

\Verse{232}
{
    \zA{黛玉 道：“不曾读书，只上了一年学，些须认得几个字。”}
}
{
    \eA{exclaimed lady Chia, sneeringly; "how could you have seen her before?"
        "Though I may not have seen her, ere this," observed Pao-yü with a smirk,
        "yet when I look at her face,}
}
{
}

\Verse{233}
{
    \zA{宝玉又道：“妹妹尊名？”}
}
{
    \eA{it seems so familiar, and to my mind, it would appear as if we had been old
        acquaintances;}
}
{
}

\Verse{234}
{
    \zA{黛玉便说了名，宝玉又道：“表字？”}
}
{
    \eA{just as if, in fact, we were now meeting after a long separation." "That
        will do! that will do!" remarked dowager lady Chia;}
}
{
}

\Verse{235}
{
    \zA{黛玉道：“无字。”}
}
{
    \eA{"such being the case, you will be the more intimate."}
}
{
}

\Verse{236}
{
    \zA{宝玉笑道：“我送妹妹 一字：莫若‘颦颦’二字极妙。”}
}
{
    \eA{Pao-yü, thereupon, went up to Tai-yü, and taking a seat next to her,
        continued to look at her again with all intentness for a good long while.
        "Have you read any books,}
}
{
}

\Verse{237}
{
    \zA{探春便道：“何处出典？”}
}
{
    \eA{cousin?" he asked. "I haven't as yet," replied Tai-yü, "read any books,}
}
{
}

\Verse{238}
{
    \zA{宝玉道：“《古今人 物通考》上说：‘西方有石名黛，可代画眉之墨。’}
}
{
    \eA{as I have only been to school for a year; all I know are simply a few
        characters." "What is your worthy name, cousin?" Pao-yü went on to ask;
        whereupon Tai-yü speedily told him her name. "Your style?" inquired Pao-yü;}
}
{
}

\Verse{239}
{
    \zA{况这妹妹眉尖若蹙，取这个字 岂不美？”}
}
{
    \eA{to which question Tai-yü replied, "I have no style." "I'll give you a
        style," suggested Pao-yü smilingly; "won't the double style 'P'in P'in,'}
}
{
}

\Verse{240}
{
    \zA{探春笑道：“只怕又是杜撰。”}
}
{
    \eA{'knitting brows,' do very well?" "From what part of the standard books does
        that come?"}
}
{
}

\Verse{241}
{
    \zA{宝玉笑道：“除了《四书》，杜撰的也 太多呢。”}
}
{
    \eA{T'an Ch'un hastily interposed. "It is stated in the Thorough Research into
        the state of Creation from remote ages to the present day," Pao-yü went on
        to explain,}
}
{
}

\Verse{242}
{
    \zA{因又问黛玉：“可有玉没有？”}
}
{
    \eA{"that, in the western quarter, there exists a stone, called Tai, (black,)
        which can be used,}
}
{
}

\Verse{243}
{
    \zA{众人都不解。}
}
{
    \eA{in lieu of ink, to blacken the eyebrows with.}
}
{
}

\Verse{244}
{
    \zA{黛玉便忖度着：“因他有 玉，所以才问我的。”}
}
{
    \eA{Besides the eyebrows of this cousin taper in a way, as if they were
        contracted, so that the selection of these two characters is most
        appropriate,}
}
{
}

\Verse{245}
{
    \zA{便答道：“我没有玉。}
}
{
    \eA{isn't it?" "This is just another plagiarism, I fear," observed T'an Ch'un,}
}
{
}

\Verse{246}
{
    \zA{你那玉也是件稀罕物儿，岂能人人皆 有？”}
}
{
    \eA{with an ironic smirk. "Exclusive of the Four Books," Pao-yü remarked
        smilingly, "the majority of works are plagiarised; and is it only I,
        perchance,}
}
{
}

\Verse{247}
{
    \zA{宝玉听了，登时发作起狂病来，摘下那玉就狠命摔去，骂道：“什么罕物! 人的高下不识，还说灵不灵呢!我也不要这劳什子！”}
}
{
    \eA{who plagiarise? Have you got any jade or not?" he went on to inquire,
        addressing Tai-yü, (to the discomfiture) of all who could not make out what
        he meant. "It's because he has a jade himself," Tai-yü forthwith reasoned
        within her mind, "that he asks me whether I have one or not.--No;}
}
{
}

\Verse{248}
{
    \zA{吓的地下众人一拥争去拾玉。}
}
{
    \eA{I haven't one," she replied. "That jade of yours is besides a rare object,
        and how could every one have one?"}
}
{
}

\Verse{249}
{
    \zA{贾母急的搂了宝玉道“孽障!你生气要打骂人容易，何苦摔那命根子！”}
}
{
    \eA{As soon as Pao-yü heard this remark, he at once burst out in a fit of his
        raving complaint, and unclasping the gem, he dashed it disdainfully on the
        floor. "Rare object, indeed!" he shouted, as he heaped invective on it;}
}
{
}

\Verse{250}
{
    \zA{宝玉满面 泪痕哭道：“家里姐姐妹妹都没有，单我有，我说没趣儿；如今来了这个神仙似的 妹妹也没有，可知这不是个好东西。”}
}
{
    \eA{"it has no idea how to discriminate the excellent from the mean, among human
        beings; and do tell me, has it any perception or not? I too can do without
        this rubbish!" All those, who stood below, were startled; and in a body they
        pressed forward, vying with each other as to who should pick up the gem.
        Dowager lady Chia was so distressed that she clasped Pao-yü in her embrace.}
}
{
}

\Verse{251}
{
    \zA{贾母忙哄他道：“你这妹妹原有玉来着。}
}
{
    \eA{"You child of wrath," she exclaimed. "When you get into a passion, it's easy
        enough for you to beat and abuse people; but what makes you fling away that
        stem of life?"}
}
{
}

\Verse{252}
{
    \zA{因 你姑妈去世时，舍不得你妹妹，无法可处，遂将他的玉带了去，一则全殉葬之礼， 尽你妹妹的孝心；二则你姑妈的阴灵儿也可权作见了你妹妹了。}
}
{
    \eA{Pao-yü's face was covered with the traces of tears. "All my cousins here,
        senior as well as junior," he rejoined, as he sobbed, "have no gem, and if
        it's only I to have one, there's no fun in it, I maintain! and now comes
        this angelic sort of cousin, and she too has none, so that it's clear enough
        that it is no profitable thing."}
}
{
}

\Verse{253}
{
    \zA{因此他说没有，也 是不便自己夸张的意思啊。}
}
{
    \eA{Dowager lady Chia hastened to coax him. "This cousin of yours," she
        explained, "would, under former circumstances, have come here with a jade;}
}
{
}

\Verse{254}
{
    \zA{你还不好生带上，仔细你娘知道！”}
}
{
    \eA{and it's because your aunt felt unable, as she lay on her death-bed, to
        reconcile herself to the separation from your cousin,}
}
{
}

\Verse{255}
{
    \zA{说着便向丫鬟手中 接来亲与他带上。}
}
{
    \eA{that in the absence of any remedy, she forthwith took the gem belonging to
        her (daughter), along with her (in the grave);}
}
{
}

\Verse{256}
{
    \zA{宝玉听如此说，想了一想，也就不生别论。}
}
{
    \eA{so that, in the first place, by the fulfilment of the rites of burying the
        living with the dead might be accomplished the filial piety of your cousin;
        and in the second place,}
}
{
}

\Verse{257}
{
    \zA{当下奶娘来问黛玉房舍，贾母便说：“将宝玉挪出来，同我在套间暖阁里，把 你林姑娘暂且安置在碧纱厨里。}
}
{
    \eA{that the spirit of your aunt might also, for the time being, use it to
        gratify the wish of gazing on your cousin. That's why she simply told you
        that she had no jade; for she couldn't very well have had any desire to give
        vent to self-praise. Now, how can you ever compare yourself with her? and
        don't you yet carefully and circumspectly put it on?}
}
{
}

\Verse{258}
{
    \zA{等过了残冬，春天再给他们收拾房屋，另作一番安 置罢。”}
}
{
    \eA{Mind, your mother may come to know what you have done!" As she uttered these
        words, she speedily took the jade over from the hand of the waiting-maid,
        and she herself fastened it on for him.}
}
{
}

\Verse{259}
{
    \zA{宝玉道：“好祖宗，我就在碧纱厨外的床上很妥当。}
}
{
    \eA{When Pao-yü heard this explanation, he indulged in reflection, but could not
        even then advance any further arguments. A nurse came at the moment and
        inquired about Tai-yü's quarters,}
}
{
}

\Verse{260}
{
    \zA{又何必出来，闹的老 祖宗不得安静呢？”}
}
{
    \eA{and dowager lady Chia at once added, "Shift Pao-yü along with me, into the
        warm room of my suite of apartments, and put your mistress, Miss Lin,}
}
{
}

\Verse{261}
{
    \zA{贾母想一想说：“也罢了。”}
}
{
    \eA{temporarily in the green gauze house; and when the rest of the winter is
        over,}
}
{
}

\Verse{262}
{
    \zA{每人一个奶娘并一个丫头照管， 馀者在外间上夜听唤。}
}
{
    \eA{and repairs are taken in hand in spring in their rooms, an additional wing
        can be put up for her to take up her quarters in." "My dear ancestor,"
        ventured Pao-yü; "the bed I occupy outside the green gauze house is very
        comfortable;}
}
{
}

\Verse{263}
{
    \zA{一面早有熙凤命人送了一顶藕合色花帐并锦被缎褥之类。}
}
{
    \eA{and what need is there again for me to leave it and come and disturb your
        old ladyship's peace and quiet?" "Well, all right," observed dowager lady
        Chia, after some consideration; "but let each one of you have a nurse, as
        well as a waiting-maid to attend on you;}
}
{
}

\Verse{264}
{
    \zA{黛 玉只带了两个人来，一个是自己的奶娘王嬷嬷，一个是十岁的小丫头，名唤雪雁。}
}
{
    \eA{the other servants can remain in the outside rooms and keep night watch and
        be ready to answer any call." At an early hour, besides, Hsi-feng had sent a
        servant round with a grey flowered curtain, embroidered coverlets and satin
        quilts and other such articles. Tai-yü had brought along with her only two
        servants; the one was her own nurse,}
}
{
}

\Verse{265}
{
    \zA{贾母见雪雁甚小，一团孩气，王嬷嬷又极老，料黛玉皆不遂心，将自己身边一个二 等小丫头名唤鹦哥的与了黛玉。}
}
{
    \eA{dame Wang, and the other was a young waiting-maid of sixteen, whose name was
        called Hsüeh Yen. Dowager lady Chia, perceiving that Hsüeh Yen was too
        youthful and quite a child in her manner, while nurse Wang was, on the other
        hand, too aged, conjectured that Tai-yü would, in all her wants, not have
        things as she liked, so she detached two waiting-maids,}
}
{
}

\Verse{266}
{
    \zA{亦如迎春等一般，每人除自幼乳母外，另有四个教 引嬷嬷，除贴身掌管钗钏盥沐两个丫头外，另有四五个洒扫房屋来往使役的小丫 头。}
}
{
    \eA{who were her own personal attendants, named Tzu Chüan and Ying Ko, and
        attached them to Tai-yü's service. Just as had Ying Ch'un and the other
        girls, each one of whom had besides the wet nurses of their youth, four
        other nurses to advise and direct them, and exclusive of two personal maids
        to look after their dress and toilette,}
}
{
}

\Verse{267}
{
    \zA{当下王嬷嬷与鹦哥陪侍黛玉在碧纱厨内，宝玉乳母李嬷嬷并大丫头名唤袭人的 陪侍在外面大床上。}
}
{
    \eA{four or five additional young maids to do the washing and sweeping of the
        rooms and the running about backwards and forwards on errands. Nurse Wang,
        Tzu Chüan and other girls entered at once upon their attendance on Tai-yü in
        the green gauze rooms, while Pao-yü's wet-nurse, dame Li, together with an
        elderly waiting-maid, called Hsi Jen, were on duty in the room with the
        large bed.}
}
{
}

\Verse{268}
{
    \zA{原来这袭人亦是贾母之婢，本名蕊珠，贾母因溺爱宝玉，恐宝 玉之婢不中使，素喜蕊珠心地纯良，遂与宝玉。}
}
{
    \eA{This Hsi Jen had also been, originally, one of dowager lady Chia's
        servant-girls. Her name was in days gone by, Chen Chu. As her venerable
        ladyship, in her tender love for Pao-yü, had feared that Pao-yü's servant
        girls were not equal to their duties, she readily handed her to Pao-yü,}
}
{
}

\Verse{269}
{
    \zA{宝玉因知他本姓花，又曾见旧人诗 句有“花气袭人”之句，遂回明贾母，即把蕊珠更名袭人。}
}
{
    \eA{as she had hitherto had experience of how sincere and considerate she was at
        heart. Pao-yü, knowing that her surname was at one time Hua, and having once
        seen in some verses of an ancient poet, the line "the fragrance of flowers
        wafts itself into man," lost no time in explaining the fact to dowager lady
        Chia, who at once changed her name into Hsi Jen.}
}
{
}

\Verse{270}
{
    \zA{却说袭人倒有些痴处：伏侍贾母时，心中只有贾母；如今跟了宝玉，心中又只 有宝玉了。}
}
{
    \eA{This Hsi Jen had several simple traits. While in attendance upon dowager
        lady Chia, in her heart and her eyes there was no one but her venerable
        ladyship, and her alone; and now in her attendance upon Pao-yü, her heart
        and her eyes were again full of Pao-yü,}
}
{
}

\Verse{271}
{
    \zA{只因宝玉性情乖僻，每每规谏，见宝玉不听，心中着实忧郁。}
}
{
    \eA{and him alone. But as Pao-yü was of a perverse temperament and did not heed
        her repeated injunctions, she felt at heart exceedingly grieved. At night,
        after nurse Li had fallen asleep,}
}
{
}

\Verse{272}
{
    \zA{是晚宝玉 李嬷嬷已睡了，他见里面黛玉鹦哥犹未安歇，他自卸了妆，悄悄的进来，笑问：“姑 娘怎么还不安歇？”}
}
{
    \eA{seeing that in the inner chambers, Tai-yü, Ying Ko and the others had not as
        yet retired to rest, she disrobed herself, and with gentle step walked in.
        "How is it, miss," she inquired smiling, "that you have not turned in as
        yet?"}
}
{
}

\Verse{273}
{
    \zA{黛玉忙笑让：“姐姐请坐。”}
}
{
    \eA{Tai-yü at once put on a smile. "Sit down, sister," she rejoined, pressing
        her to take a seat.}
}
{
}

\Verse{274}
{
    \zA{袭人在床沿上坐了。}
}
{
    \eA{Hsi Jen sat on the edge of the bed. "Miss Lin," interposed Ying Ko
        smirkingly,}
}
{
}

\Verse{275}
{
    \zA{鹦哥笑道： “林姑娘在这里伤心，自己淌眼抹泪的，说：‘今儿才来了，就惹出你们哥儿的病 来。}
}
{
    \eA{"has been here in an awful state of mind! She has cried so to herself, that
        her eyes were flooded, as soon as she dried her tears. 'It's only to-day
        that I've come,' she said, 'and I've already been the cause of the outbreak
        of your young master's failing.}
}
{
}

\Verse{276}
{
    \zA{倘或摔坏了那玉，岂不是因我之过！’}
}
{
    \eA{Now had he broken that jade, as he hurled it on the ground, wouldn't it have
        been my fault? Hence it was that she was so wounded at heart,}
}
{
}

\Verse{277}
{
    \zA{所以伤心，我好容易劝好了。”}
}
{
    \eA{that I had all the trouble in the world, before I could appease her."
        "Desist at once, Miss! Don't go on like this,"}
}
{
}

\Verse{278}
{
    \zA{袭人道： “姑娘快别这么着!将来只怕比这更奇怪的笑话儿还有呢。}
}
{
    \eA{Hsi Jen advised her; "there will, I fear, in the future, happen things far
        more strange and ridiculous than this; and if you allow yourself to be
        wounded and affected to such a degree by a conduct such as his,}
}
{
}

\Verse{279}
{
    \zA{若为他这种行状你多心 伤感，只怕你还伤感不了呢。}
}
{
    \eA{you will, I apprehend, suffer endless wounds and anguish; so be quick and
        dispel this over-sensitive nature!" "What you sisters advise me,"}
}
{
}

\Verse{280}
{
    \zA{快别多心。”}
}
{
    \eA{replied Tai-yü, "I shall bear in mind,}
}
{
}

\Verse{281}
{
    \zA{黛玉道：“姐姐们说的，我记着就是了。”}
}
{
    \eA{and it will be all right." They had another chat, which lasted for some
        time, before they at length retired to rest for the night.}
}
{
}

\Verse{282}
{
    \zA{又叙了一回，方才安歇。}
}
{
    \eA{The next day, (she and her cousins) got up at an early hour and went over to
        pay their respects to dowager lady Chia,}
}
{
}

\Verse{283}
{
    \zA{次早起来省过贾母，因往王夫人处来。}
}
{
    \eA{after which upon coming to madame Wang's apartments, they happened to find
        madame Wang and Hsi-feng together,}
}
{
}

\Verse{284}
{
    \zA{正值王夫人与熙凤在一处拆金陵来的书 信，又有王夫人的兄嫂处遣来的两个媳妇儿来说话。}
}
{
    \eA{opening the letters which had arrived from Chin Ling. There were also in the
        room two married women, who had been sent from madame Wang's elder brother's
        wife's house to deliver a message. Tai-yü was, it is true, not aware of what
        was up,}
}
{
}

\Verse{285}
{
    \zA{黛玉虽不知原委，探春等却晓 得是议论金陵城中居住的薛家姨母之子――表兄薛蟠，倚财仗势，打死人命，现在 应天府案下审理。}
}
{
    \eA{but T'an Ch'un and the others knew that they were discussing the son of her
        mother's sister, married in the Hsüeh family, in the city of Chin Ling, a
        cousin of theirs, Hsüeh P'an, who relying upon his wealth and influence had,}
}
{
}

\Verse{286}
{
    \zA{如今舅舅王子腾得了信，遣人来告诉这边，意欲唤取进京之意。}
}
{
    \eA{by assaulting a man, committed homicide, and who was now to be tried in the
        court of the Ying T'ien Prefecture. Her maternal uncle, Wang Tzu-t'eng, had
        now, on the receipt of the tidings,}
}
{
}

\Verse{287}
{
    \zA{毕竟怎的，下回分解。}
}
{
    \eA{despatched messengers to bring over the news to the Chia family.}
}
{
}

\Verse{288}
{
    \zA{(本章完)}
}
{
    \eA{But the next chapter will explain what was the ultimate issue of the wish
        entertained in this mansion to send for the Hsüeh family to come to the
        capital.}
}
{
}
